\documentclass[11pt, a4paper]{article}

% PACKAGES
\usepackage[utf8]{inputenc}
\usepackage{amsmath}
\usepackage{amssymb}
\usepackage[margin=1in]{geometry}
\usepackage{authblk}


\title{A Reflection on the Localization of Matter, Elementary Particles, and Their Interactions}
\author{Haifeng Qiu}
%\date{}

\begin{document}

\maketitle

\begin{abstract}
In this paper, we continue to develop an in-depth and open-ended reflection on ``matter,'' the most fundamental physical reality, within the theoretical framework of ``Computational Realism'' that we are constructing. We hope readers will view this article as a candid ``research log.'' Our purpose is to invite you to embark on this journey of logical reasoning with us. In the process, we shall discover that some seemingly simple principles may, in an almost inevitable manner, lead to a series of interesting and profound phenomena and conclusions.

We do not claim that the model proposed here is the final answer. We are acutely aware that any narrative attempting to reconstruct the entirety of particle physics from its foundations will inevitably face the most rigorous scrutiny and may even be classified as a non-mainstream exploration. We accept this with equanimity. The sole purpose of sharing these thoughts is the hope that this unified and logically self-consistent perspective can offer you—our respected fellow explorers—some new and perhaps beneficial angles of thought. If any deductive step within this text brings you a moment of intellectual resonance regarding the nature of reality, then the value of our work has found its highest affirmation.
\end{abstract}

\section{The Localization Principle---Emergence from Chaos}

Within the framework of ``Computational Realism,'' all fundamental existences—whether matter as a ``structure'' (soliton) or light as a ``process'' (photon)—owe their ability to exist as stable, localized entities to the same profound dynamical principle.

We name this the \textbf{Minimal Spatial Encoding Principle}. This principle asserts that any \textbf{Coherent Information Pattern} propagating through a discrete, chaotic computational field is constrained by two synergistic dynamical mechanisms, causing it to spontaneously and non-linearly localize within a finite spatiotemporal scale.

To analyze this principle, we divide an information pattern into \textbf{the whole} and \textbf{the part}, whose dynamics are governed by the following two fundamental hypotheses:

\subsection{Mechanism One: The Cohesive Effect of the Self-Generated Field}

\begin{itemize}
    \item \textit{Hypothesis 1 (Interference / Time Dilation Effect):} The propagation of periodic information slows down its effective spatial velocity in the vicinity of the holistic, same-frequency information due to an \textbf{interference effect} or \textbf{“bandwidth occupation.”} Physically, this effect is analogous to the self-focusing of the \textbf{Kerr Effect}.
    
    \item \textbf{Inference:} According to the principle of least action, this ``slow medium'' effect, generated by the pattern itself, creates a powerful tendency for \textbf{Self-focusing}. This mechanism is the primary contributor to the pattern's \textbf{Cohesion}, causing its information to converge toward the center.
\end{itemize}

\subsection{Mechanism Two: The Boundary Effect of the Chaotic Background}

\begin{itemize}
    \item \textit{Hypothesis 2 (Random Reflection Effect):} As information propagates through a discrete, random computational field, its outwardly dispersing ``wavefront'' undergoes complex interactions with countless random ``bit defects'' (manifested by the `$b_0$' chaos) and is thereby \textbf{probabilistically reflected} back toward the center. The physical essence of this mechanism is homologous to the phenomenon of \textbf{Anderson Localization}.
    
    \item \textbf{Inference:} This effect provides the pattern with an effective \textbf{``Probabilistic Boundary'' or ``wall,''} ensuring that its information does not dissipate and escape indefinitely.
\end{itemize}

\subsection{Combined Effect and the Emergence of Physical Reality}

These two mechanisms work in synergy: Mechanism One (self-focusing) provides a powerful \textbf{cohesive force} from the \textbf{inside}, causing the pattern to tend to aggregate; Mechanism Two (chaotic reflection) provides a robust \textbf{constraining force} from the \textbf{outside}, preventing the pattern from dispersing infinitely.

A stable material particle (soliton) is precisely the non-linear resonant structure formed when the tendency of information to diffuse and the effect of its cohesion achieve a perfect dynamic equilibrium.

The \textbf{Directionality} of the pattern's propagation originates from a definite, quantized \textbf{``focused initial momentum''} imparted by the geometric asymmetry of the source during its creation event (e.g., an atomic transition).

Profoundly, \textbf{Mechanism One (self-focusing/time dilation effect)} is the same fundamental mechanism as \textbf{gravity} in our theory; gravity is the macroscopic average manifestation of this effect across all computational layers. \textbf{Macroscopic matter}, due to its immense average effect, is predominantly governed by Mechanism One and is almost unaffected by the microscopic effects of Mechanism Two (random reflection), thus allowing it to exist in dispersed forms. The prominence of these two effects is a \textbf{Scale-Dependent} issue.

\section{The Division of Forces and the Constitution of Matter---From Gradient Force to Binding Force}

\subsection{The Unified and Hierarchical Definition of ``Charge'': An ``Information Fountain'' Model}

\begin{itemize}
    \item \textbf{Unified Physical Picture of Charge:}  
    At the center of an information soliton (such as an electron or a quark), there exists a continuous, holistic \textbf{cross-computational-layer information flow} (either upward or downward). This core information flow, on \textbf{every} computational layer, \textbf{expands outward} in the spatial dimensions to form an annular flow field. Then, beyond a certain spatial distance, it \textbf{re-converges and sinks/rises} back to an adjacent layer, forming a self-consistent, closed \textbf{``Information Torus''} or \textbf{``Information Fountain.''} We must emphasize that this "Information Torus" is a simplified description in a 2D context for conceptual clarity; the physical reality is a "3D Information Fountain" unfolding in three-dimensional space.
    
    What we call \textbf{``charge'' is the macroscopic manifestation of this ``Information Fountain'' at different computational layers}.
    
    \item \textbf{Hierarchical Definition of Charges:}
    \begin{itemize}
        \item \textbf{Electric Charge:} Is the \textbf{torus} formed by this ``Information Fountain'' at the \textbf{`$b_0-b_1$'} layer. This torus can extend over a \textbf{certain spatial distance}, forming a spatial gradient field.
    
        \item \textbf{Strong Charge:} Is the \textbf{torus} formed by this ``Information Fountain'' at the \textbf{`$b_1-b_2$'} layer. Due to the constraint of the ``axiom of parasitic existence'' described below, this higher-level torus \textbf{cannot} exist over long distances detached from its `$b_0-b_1$' energy base. It is therefore \textbf{confined} within an \textbf{extremely short spatial region}.
    \end{itemize}
    
    \item \textbf{Cross-Layer Transport Protocols:}
    \begin{itemize}
        \item \textbf{Unidirectionality Constraint:} Within a stable soliton, the overall flow direction of this ``Information Fountain'' must be uniform (either a holistic upward ``fountain'' or a holistic downward ``funnel'').
        \item \textbf{Bandwidth Quantization:} The minimal information unit for the `$b_0-b_1$' channel is \textbf{1}; the minimal information unit for the `$b_1-b_2$' channel is \textbf{1/3}.
    \end{itemize}
\end{itemize}

\subsection{Electromagnetic Interaction: An ``Autonomous Gradient Force''}

\begin{itemize}
    \item \textbf{Definition of the Electron:} The electron is a simple, stable information soliton involving only the information flow of the `$b_0-b_1$' layer.
    \item \textbf{Nature of Electromagnetic Force:} A charged soliton (like an electron) maintains a static information gradient field (electric field) on the surrounding `$b_0-b_1$' computational layers. The electromagnetic force is the interaction of a charged soliton with the gradient field produced by other charges. The gradient field produced by a single charge is maintained only over a finite spatial distance; the macroscopic electric field is the average effect of a large number of charges.
\end{itemize}

\subsection{Strong Interaction: A ``Parasitic Binding Force''}

\begin{itemize}
    \item \textbf{Definition of the Quark:} A quark is a composite information soliton spanning three layers, `$b_0-b_1-b_2$', which carries both \textbf{electric charge} and \textbf{strong charge}.
    
    \item \textbf{The Fundamental Physics of Strong Charge: From ``Parasitic Existence'' to ``Rotational Symmetry''}
    \begin{enumerate}
        \item \textbf{A Foundational Axiom—The Parasitic Existence of Strong Charge:} The existence of the strong charge field (the information flow on the `$b_1-b_2$' layer) requires an active `$b_0-b_1$' information flow to serve as its \textbf{dynamical foundation}.
        \item \textbf{Inference 1 (Short-Range):} As a direct consequence of this axiom, the strong charge field is a \textbf{short-range field} enforced by dynamical laws.
        \item \textbf{Inference 2 (Rotational Effect):} For a field that cannot diffuse, its only stable dynamical form is to \textbf{internalize} all its ``computational activity'' into a local, closed \textbf{``internal rotation'' or ``phase.''}
        \item \textbf{Inference 3 (Rotational Group Symmetry):} The mathematical language describing this ``internal phase rotation'' \textbf{must be a rotational group}. Based on the observed rules of hadron composition, we find that the properties of this group perfectly match the \textbf{`$U(1)$' rotation group}, whose algebra is equivalent to \textbf{“mod 1 arithmetic.”} Under this algebra, any integer-valued net rotation is equivalent to zero (`$n \pmod{1} = 0$').
    \end{enumerate}
    
    \item \textbf{The Law of Hadron Formation:}
    \begin{itemize}
        \item \textbf{Axiom of Quantized Existence:} Any ``information pattern'' that can exist stably as an independent, self-consistent entity (be it a lepton or a hadron) must have a total \textbf{electric charge} that is an \textbf{integer multiple} of the fundamental unit.
        \item \textbf{Charge Conservation Constraint:} For any stable quark pattern, the \textbf{absolute value of the sum of its electric charge and strong charge} must equal the fundamental information unit of the `$b_0-b_1$' channel, which is \textbf{1}. $|$Electric Charge + Strong Charge$| = 1$.
        \item \textbf{Quark Charge Assignments:} Based on the above axioms and constraints, only four stable quark patterns exist in the first generation:
        \begin{enumerate}
            \item \textbf{Down Quark (d):} Electric Charge `-1/3', Strong Charge `-2/3'
            \item \textbf{Up Quark (u):} Electric Charge `+2/3', Strong Charge `+1/3'
            \item \textbf{Anti-down Quark ($\bar{d}$):} Electric Charge `+1/3', Strong Charge `+2/3'
            \item \textbf{Anti-up Quark ($\bar{u}$):} Electric Charge `-2/3', Strong Charge `-1/3'
        \end{enumerate}
        \item \textbf{Emergence of Strong-Charge Neutrality:} As a necessary \textbf{theorem}, any hadron that satisfies the ``Axiom of Quantized Existence'' (i.e., has an integer total electric charge) must have a \textbf{total strong charge} of all its constituent quarks that is \textbf{zero} in the sense of `$U(1)$'. (For example, the total strong charge of a proton `uud' is `$(+1/3) + (+1/3) + (-2/3) = 0 \pmod{1}$').
    \end{itemize}
    
    \item \textbf{Nature of the Nuclear Force:}
    \begin{itemize}
        \item An isolated nucleon is \textbf{strong-charge neutral}.
        \item However, when two nucleons come extremely close, their respective internal \textbf{information circuits overlap}. This causes the ``strong charge'' information flow, which was perfectly circulating within each nucleon, to \textbf{``leak,'''} forming a new, \textbf{short-range, inter-nucleon information flow circuit}. The powerful attractive force generated by this newly formed circuit is the fundamental source of the ``nuclear force.'' This mechanism is dynamically isomorphic to the \textbf{van der Waals force} between molecules.
    \end{itemize}
    
    \item \textbf{Quark Confinement:}
    \begin{itemize}
        \item \textbf{Intrinsic Reason:} The fundamental reason for confinement is that all quarks within a stable hadron (baryon or meson) must be \textbf{``parasitic'' upon the same fundamental information flow structure of the `$b_0-b_1$' layer}. An isolated quark has an incomplete `$b_0-b_1$' base structure, violating the Axiom of Quantized Existence, and thus cannot exist stably on its own.
        \item \textbf{Dynamical Manifestation:} The ``strong charge force'' between quarks is extraordinarily large. Any attempt to separate a quark from a strong-charge neutral nucleon requires an immense injection of energy. Due to the extremely small scale of quarks, any probe capable of acting on them will inevitably inject energy into the system \textbf{far exceeding the creation threshold}. The final result is not ``separation'' but the creation of new quark pairs from the vacuum, forming new, stable hadrons that satisfy all conservation laws and neutrality requirements.
    \end{itemize}
\end{itemize}

\section{Weak Interaction---The ``Decay'' and ``Transmutation'' of Solitons}

In the framework of ``Computational Realism,'' the ontological status of the weak interaction is fundamentally reconstructed. It is no longer considered an ``interaction force'' parallel to the other three but is understood as an \textbf{intrinsic, dynamical instability} of the ``information soliton'' itself.

\subsection{The Essence of Weak Interaction: A Process of ``Topological Reconfiguration''}

\begin{itemize}
    \item \textbf{Definition:} The weak interaction describes the dynamical process wherein a metastable ``information soliton'' (e.g., a neutron), under the combined influence of \textbf{internal factors (intrinsic structural imperfections)} and \textbf{external factors (continuous energy/information exchange with the environment, such as `$b_0$' chaotic fluctuations or high-energy collisions)}, is triggered to undergo a \textbf{Holistic Topological Reconfiguration}.
    \item \textbf{Core Mechanism:} This is a process where internal bit information \textbf{``collapses'' and ``ruptures''} from the ordered soliton structure, with the ultimate goal of allowing the system to evolve into a new combination of solitons with lower total energy and greater structural stability. It is fundamentally a \textbf{Process}, not an instantaneous event.
\end{itemize}

\subsection{Reinterpretation of the ``Weak Force'': A ``Driving Force for the Diffusion'' of Internal Information}

\begin{itemize}
    \item \textbf{Definition:} In our model, the ``weak force'' is not a ``force'' in the traditional sense describing interactions between solitons. It should be understood more as the \textbf{driving force for the diffusion} of internal information (energy and momentum) that was ``confined'' by the soliton's structure during that ``topological reconfiguration'' process.
    \item \textbf{Physical Picture:} It is like a high-pressure vessel (the metastable soliton) suddenly rupturing, and its high-pressure gas (the confined information) violently erupts outward. The driving force of this ``eruption'' is the dynamical essence of the ``weak force.''
\end{itemize}

\subsection{The Essence of W/Z Bosons: High-Energy ``Intermediate Resonant States''}

\begin{itemize}
    \item \textbf{Definition:} The W/Z bosons are \textbf{real, yet ephemeral, ``High-Energy Intermediate Resonances''} formed by the immense internal energy that is temporarily ``deconfined'' during the violent ``topological reconfiguration'' event.
    \item \textbf{Reinterpretation of Mass:} A value such as `$80.4~\text{GeV}$' can be understood as an \textbf{``instantaneous effective rest mass.''} It represents a fragment in the process of rupture, a \textbf{``temporary information packet,''} and signifies the \textbf{upper limit of energy it can stably bind} before it completely disintegrates. Any energy injected beyond this limit will exist as kinetic energy or cause it to split into separate ``information packets.''
\end{itemize}

\subsection{The Core ``Signatures'' of Weak Interaction}

\begin{itemize}
    \item \textbf{Flavor Changing \& CKM Matrix:}
    \begin{itemize}
        \item \textbf{Phenomenon:} The weak interaction is the only force that can change the ``flavor'' (generation) of quarks and leptons, and the probability of these cross-generational transitions is described by a complex CKM matrix.
        \item \textbf{Theoretical Explanation:}
        \begin{itemize}
            \item \textbf{Essence of ``Flavor'':} Different ``flavors'' (e.g., u vs. s quark) correspond to solitons with \textbf{different topological structures of ``cross-layer information flow.''}
            \item \textbf{Changing ``Flavor'':} The ``weak interaction,'' being a violent ``topological reconfiguration,'' naturally includes the possibility of ``rewiring'' these cross-layer information flows, making it the only mechanism that can change ``flavor.''
            \item \textbf{Origin of the CKM Matrix:} The values in the CKM matrix represent the \textbf{``transition probability amplitudes''} for jumping from an initial topological configuration (e.g., an s quark) to another possible final state topology (e.g., a u quark) during the complex dynamical process of ``topological reconfiguration.'' These values reflect the \textbf{``convertibility'' or ``similarity''} between different topological structures.
        \end{itemize}
    \end{itemize}
\end{itemize}

\section{A Unifying Conjecture---The Common Origin of Neutrinos, Parity Violation, and Dark Matter}

Having established our dynamical model for the localization of matter and fundamental interactions, we find that the internal logic of this framework leads us, almost inevitably, to a unified, albeit conjectural, yet logically highly self-consistent explanation for three of the most profound mysteries in the universe: the nature of neutrinos, the origin of parity violation, and the existence of dark matter.

This chapter aims to share this exploratory line of thought. We do not claim it to be the final answer, but rather hope to demonstrate how, by strictly adhering to the core axiom of ``hierarchical computation,'' an unexpectedly new cosmological picture may unfold before us.

\subsection{The Neutrino Puzzle: A Messenger from a ``Parallel Family''?}

\noindent\textbf{The Observed Puzzle:}\\
Neutrinos exhibit a series of extremely unique, ``anomalous'' properties: 1) They are nearly ``ghostly,'' participating only in weak and gravitational interactions, while completely avoiding electromagnetic interactions (electrically neutral); 2) Their mass is extremely small, far below that of all other matter particles. Why does such a particle, almost ``insulated'' from our visible world, exist?

\noindent\textbf{A Possible Solution (Our Conjecture):}\\
We conjecture that the uniqueness of the neutrino stems from the fact that it does not belong to the ``visible matter family'' we are familiar with at all.

\begin{itemize}
    \item \textbf{The ``Visible Family'' Initiated at the `$b_0$' Layer:} We propose that all matter in our universe perceivable by the electromagnetic force (electrons, quarks) shares a common feature: their existence is \textbf{``constructed upon an initiation at the `$b_0$' layer.''} Their stable structures are profoundly dependent on direct interaction with the bottom-most `$b_0$' chaos, the macroscopic manifestation of which is ``electric charge.''

    \item \textbf{Neutrinos—The ``Dimly Lit Family'' Initiated at the `$b_1$' Layer:} In contrast, we conjecture that the neutrino family is the first and simplest particle family initiated purely at the \textbf{`$b_1$' layer}. This conjecture provides a unified dynamical explanation for the core puzzles of the neutrino:
    \begin{enumerate}
        \item \textbf{Why is it electrically neutral?} \\
        Because the core stable structure of the neutrino \textbf{lacks a `$b_0-b_1$' cross-layer information flow loop}. The influence of its `$b_1-b_2$' loop on the `$b_0-b_1$' layer is extremely weak, possibly because their \textbf{resonant frequencies are too different}; meanwhile, on a macroscopic scale, the \textbf{densities of neutrinos and anti-neutrinos in the cosmic background are equal}, and their faint secondary effects would statistically cancel each other out.

        \item \textbf{Why is its mass so small?} \\
		The rest mass of a particle primarily originates from the structural cost required to overcome the \textbf{size ``detuning'' of the ``computational cores''} of different computational layers during its cross-layer construction. The neutrino, as an independent soliton of the `$b_1-b_2$' layer, has a smaller ``detuning'' cost compared to a `$b_0-b_1$' layer soliton. It therefore requires a smaller coupling structure to exist stably, resulting in a small rest mass.
    \end{enumerate}
\end{itemize}

\subsection{The Parity Violation Puzzle: A Dynamic, Local, and Relativistic Effect}

\noindent\textbf{The Observed Puzzle:}\\
The weak interaction—the force intimately linked with neutrinos—maximally violates ``mirror symmetry'' (parity). Experiments have confirmed that, at the level of the weak interaction, the universe is staunchly ``left-handed.'' What is the origin of this profound, fundamental asymmetry?

\noindent\textbf{A Possible Solution (Our Conjecture):}\\
We conjecture that parity violation is not a global, fixed ``rule of the universe,'' but rather a \textbf{dynamic, local, and necessary emergence from the core axiom of ``hierarchical computation'' in our theory}. Its essence is a profound \textbf{``Principle of Hierarchical Relativity.''}

\begin{itemize}
    \item \textbf{Principle of Hierarchical Relativity:} We hypothesize that the `$b_0$' and `$b_1$' layers are two \textbf{distinct, co-equal dynamical frames of reference}. Between them, within every region of space, there exists an \textbf{ever-evolving, intrinsic ``relative rotation.''} Chirality is a measure of a particle's proximity to the rotation of the `$b_0$' layer versus the `$b_1$' layer.

    \item \textbf{The Emergence of Symmetry and Asymmetry:} \\
    Whether a particle exhibits chiral symmetry depends entirely on its own \textbf{``computational layer composition''} and its coupling relationship with these two different dynamical frames.
    \begin{enumerate}
        \item \textbf{Visible Matter (`$b_0$' Family)—A Chiral-Symmetric ``Mixed State'':}\\
        A ``normal,'' visible particle (like an electron), being an entity \textbf{initiated at the `$b_0$' layer}, exists stably as a \textbf{``mixed state''} spanning both `$b_0$' and `$b_1$'. Therefore, it must possess two fundamental dynamical modes to define its spin/chirality: one mode that is \textbf{``closer to `$b_0$'''} (which we define as right-handed), and another that is \textbf{``closer to `$b_1$'''} (which we define as left-handed). Possessing these two stable states allows visible particles to appear, on the whole, as \textbf{chiral-symmetric}.

        \item \textbf{Neutrinos (`$b_1$' Family)—A Chiral-Asymmetric ``Pure State'':}\\
        The neutrino, as an entity purely \textbf{initiated at the `$b_1$' layer}, \textbf{does not have the option of a state that is ``closer to `$b_0$'''}. From the perspective of particles in the `$b_0$' family, its only stable dynamical mode is one that remains perfectly \textbf{synchronized} with the ``relative rotation'' characteristic of the `$b_1$' frame in which it resides. Therefore, from our `$b_0$' observer's perspective, \textbf{the neutrino is always bound to the `$b_1$' frame with that ``relative rotation,''} and thus exhibits only one chirality (namely, what we define as ``left-handed'').
    \end{enumerate}

    \item \textbf{The Dynamic and Local Nature of the Chiral Field:}\\
    We further deduce that the ``relative rotation'' between `$b_0$' and `$b_1$' is a local field that is \textbf{dynamically changing} in spacetime. The creation and decay of every chiral particle (like a neutrino), through a mechanism of ``action and reaction,'' contributes a ``counter-chirality'' to its surrounding `$b_0-b_1$' background field. It is a \textbf{macroscopic, spatiotemporally local, dynamic equilibrium}. And this field is very likely the same entity as the background field of entanglement we have discussed previously.
\end{itemize}

\subsection{The Necessary Existence of Dark Matter}
\noindent\textbf{A Necessary Inference from the Above Conjectures:}\\
Once we accept the idea that ``particle families can be initiated at different computational layers'' and successfully use it to explain the existence of neutrinos and the asymmetry of the weak force, a grander and more inevitable inference naturally emerges:

\begin{itemize}
    \item If there exists a ``visible world'' initiated at the `$b_0$' layer and a ``dimly lit world'' initiated at the `$b_1$' layer (neutrinos), then \textbf{there must exist even ``darker,'' more ``invisible'' particle families initiated at higher computational layers.}

    \item \textbf{The Ultimate Identity of Dark Matter:} We therefore propose our ultimate conjecture—that the \textbf{dark matter}, which constitutes the vast majority of the matter mass in our universe, is primarily composed of these \textbf{``Dark Particle Families''} independently ``initiated'' at the `$b_2$', `$b_3$', or even higher computational layers, which are almost entirely ``incompatible in communication protocol'' with our visible world. They do not participate in the electromagnetic, strong, or possibly even the weak force. The only, unavoidable connection they share with us is \textbf{gravity}.
\end{itemize}

\subsection{A Bolder Conjecture---The Gravity of the Neutrino}
We speculate that the ratio of inertial mass to gravitational mass for neutrinos may differ from that of matter initiated at the `$b_0$' layer.

The inertial mass and gravitational mass that we define, and their benchmark, the gravitational constant, are established on the basis of the \textbf{“`$b_0$' layer as the common resonant origin.”} For a particle like the neutrino, initiated at the `$b_1$' layer, the ratio of its inertial to gravitational mass might exhibit a \textbf{``significant separation''} compared to normal matter. Because the larger computational core of the `$b_1$' layer has a greater effect on the layers below, it is possible that the gravitational mass of the neutrino, under the standard definition, is greater than its inertial mass, or in other words, it possesses a larger gravitational constant.

By analogy, a well-known entity that violates this ratio of inertial to gravitational mass is the photon.

\section*{Conclusion}
At this point, we have completed an exploratory reflection based on ``Computational Realism.'' We have attempted, starting from first principles, to provide a unified and internally self-consistent explanation for the core mysteries of particle physics. If this thought process itself has brought you any beneficial inspiration, we are content.


\end{document}